\section*{Introduction}

A key-value store holds items under names, and every value in every item is drawn from one shared
domain of attributes rather than from a type of its own. A schema names types, so a type is a
subset of that one domain, and the code that moves values into it and back out is a retract of it.

Three questions about such code are separate, and conflating them is what makes a side condition
look optional: which terms can be written, which terms are well formed, and what a term means. A
grammar answers the first; the judgement of §3 to §5 answers the second, because a side condition
relates siblings rather than a term to its subterms; and §2's pair and its laws answer the third.

Soundness proves that every derivable term denotes a pair satisfying the four laws, with three
residues named exactly, and §6's second half refutes the equations a reader would expect to hold,
each with a witness.

The sections take it in turn. §1 fixes the attribute domain, in which every value a store holds is
drawn from one shared set. §2 says what a codec is: a partial retraction of that domain onto a set
of values, together with the equivalence its own reading induces. §3 gives the judgement two
indices, because a codec has two sides, the values it holds and the member it stores them in. §4
adds a second judgement for fields, where presence is a property of a name binding and never of an
attribute. §5 states the rules a grammar cannot state, because they relate siblings rather than a
term to its subterms. §6 settles which rewrites preserve the pair. §7 separates the three layers,
and closes on what the calculus leaves to the caller.
